\section{Introduction}
\label{sec:introduction}

Transformer models are built around attention and have become a dominant architecture for natural language processing since the introduction of the Transformer itself~\cite{vaswani2017attention}, BERT~\cite{devlin2019bert}, and compact variants such as DistilBERT~\cite{sanh2019distilbert}. Because attention weights are produced internally by the model and can be visualized directly, they are often treated as a natural explanation of which input tokens matter most for a prediction. In sentiment classification, for example, one might expect tokens with higher attention to correspond to the words that most strongly support the predicted label.

That intuition is appealing, but it remains controversial. Prior work has shown that attention weights do not necessarily align with other feature-importance measures or faithfully reflect how much a token influences the output~\cite{jain2019attention,serrano2019attention}. At the same time, later discussion argues that the usefulness of attention depends on what one expects from an explanation and how the comparison is performed~\cite{wiegreffe2019attention}. This debate makes attention a good target for a focused empirical study: if attention is to be used as an explanation, it should at least show consistent agreement with alternative importance signals or reliably identify tokens whose removal changes the model's decision.

In this paper, we ask a narrow question: \emph{do attention weights reflect token importance in a Transformer sentiment classifier?} We study this question with DistilBERT fine-tuned on SST-2, a standard binary sentiment task from GLUE~\cite{socher2013recursive,wang2018glue}. Our analysis compares three token-ranking signals on the full validation split: last-layer \texttt{[CLS]} attention, gradient $\times$ input, and leave-one-out (LOO) deletion. We evaluate them in two complementary ways. First, we measure agreement between rankings using Spearman correlation. Second, we test how disruptive each ranking is by deleting top-ranked tokens and measuring confidence drops and prediction flips.

The paper makes three contributions. First, it provides a compact multi-seed benchmark for comparing attention with gradient-based and perturbation-based token importance signals in a Transformer classifier. Second, it separates two questions that are often conflated: whether methods rank tokens similarly, and whether they identify tokens whose removal most affects the model. Third, it packages the study in a reproducible workflow that produces saved metrics, token rankings, deletion summaries, and paper figures from the same experiment pipeline.
